\chapter{变分法、Euler--Lagrange 与 Pontryagin 极小值原理}
\label{ch:opt-variational-pmp}

% ════════════════════════════════════════════════════════════════
%  控制部 C1a 首章（卷四《控制理论基础》最优控制子部之首）。
%  吸收 Robotics_Tutorial 40/10(变分法) + 40/20(PMP)，研究生密度。
%  与 C0 \cref{ch:optimal-control-basics} 三层密度梯度：导论→C0 做透→本章严格化。
%  右扰动主线（SE(3) 段对齐 \cref{ch:lie}）；PMP 用极大化约定，每个接口标 lambda→−lambda。
%  勘误 E1–E5 已应用（设计稿 §2，裁决清单：4 真错 1 部分）。
% ════════════════════════════════════════════════════════════════

\begin{note}
\textbf{本章定位与记号约定.}\;
本章是控制部最优控制子部的首章。导论 \cref{ch:control} 已给出 LQR/HJB/MPC 的全景与工程用法，现代控制基础 \cref{ch:optimal-control-basics} 已做透变分法入门、正则方程、Pontryagin 原理的\emph{陈述}与双积分器 Bang--Bang 图解。本章在这两层之上，把变分法到极小值原理的整条逻辑链\emph{严格化}：补齐第一变分的泛函分析地基、变分基本引理与针状变分的完整证明骨架、二阶充分条件、奇异弧分类与现代扩展。三层密度互不重复，凡用到 \cref{ch:optimal-control-basics} 的基础结论一律交叉引用指回，不再重证。

记号沿用全书：状态 $\mathbf x\in\mathbb R^n$、控制 $\mathbf u\in U\subseteq\mathbb R^m$、共态（co-state / 伴随 adjoint）$\boldsymbol\lambda$、值函数 $V(\mathbf x,t)$。代价泛函记 $J$；为避免与辛结构矩阵冲突，后者记 $\boldsymbol\Omega$。梯度 $\nabla$ 取列向量、Jacobian 取 $\partial\mathbf f/\partial\mathbf x$（沿用 \nameref{ch:notation} 与 \cref{ch:matderiv}）。\textbf{PMP 全章采用 Pontryagin 原始的极大化约定}（便于讲针状变分），与 \cref{ch:control}/\cref{ch:optimal-control-basics} 的极小化约定差一个共态符号；每个接口处显式标注换算 $\boldsymbol\lambda\to-\boldsymbol\lambda$，并复述工程铁律「反馈律恒为负反馈 $\mathbf u=-\mathbf K\mathbf x$」。
\end{note}

\section{从有限维到无穷维：泛函与第一变分}
\label{sec:pmp-functional}

\paragraph{动机：优化的对象是一整条曲线。}
有限维优化求的是使 $f(x_1,\dots,x_n)$ 取极值的点 $\mathbf x^\ast\in\mathbb R^n$。但大量物理与机器人问题要找的不是「一个点」，而是「一条曲线」：两点间哪条路径使小球下滑最快？机器人关节从 $A$ 到 $B$ 的哪条轨迹耗能最少？这些问题的优化对象是整个函数 $y(\cdot)$，需要一套「对函数求极值」的微积分——\emph{变分法}。\cref{ch:optimal-control-basics} 已以最短弧长（\cref{ex:oc-arclength}）作了直观引入；这里我们补齐其严格地基，因为 Pontryagin 极小值原理的证明正建立在「容许类、变分强弱、针状扰动」这些精细概念之上。

\begin{definition}[泛函与容许类]\label{def:pmp-functional}
设 $\mathcal V$ 是函数空间（如 $C^1[a,b]$）。映射 $J:\mathcal V\to\mathbb R$ 称为\emph{泛函}。最常见的是积分型泛函
\begin{equation}
  J[y]=\int_a^b L\big(x,y(x),y'(x)\big)\,\mathrm dx,
  \label{eq:pmp-integral-functional}
\end{equation}
其中 $L$ 称为 Lagrangian。固定端点问题的\emph{容许类}为
$\mathcal A=\{y\in C^k[a,b]:y(a)=A,\;y(b)=B\}$，
$k$ 取决于 $L$ 中最高阶导数（标准 $L(x,y,y')$ 取 $k=1$ 或 $2$）。
\end{definition}

\begin{insight}
泛函之于函数，恰如函数之于数：函数把数映射为数，泛函把函数映射为数。有限维的工具（梯度、极值条件、Hessian、Taylor 展开）将被「抬升」到无穷维。下表是贯穿全章的指南针——每当在无穷维迷失，回到对应的有限维概念即可找回方向。

\begin{center}\small
\begin{tabular}{ll}
\toprule
有限维 $\mathbb R^n$ & 无穷维（函数空间）\\
\midrule
向量 $\mathbf x=(x_1,\dots,x_n)$ & 函数 $y(\cdot)$\\
函数 $f:\mathbb R^n\to\mathbb R$ & 泛函 $J:\mathcal V\to\mathbb R$\\
梯度 $\nabla f$ & 变分导数 $\delta J/\delta y$\\
驻点 $\nabla f=\mathbf 0$ & Euler--Lagrange 方程（$\delta J=0$）\\
Hessian $\nabla^2 f\succeq 0$ & 第二变分 $\delta^2 J\ge 0$\\
$f(\mathbf x+\mathbf h)=f+\nabla f^\top\mathbf h+\tfrac12\mathbf h^\top\nabla^2 f\,\mathbf h+\cdots$ & $J[y+\varepsilon h]=J[y]+\varepsilon\,\delta J+\tfrac{\varepsilon^2}{2}\delta^2 J+\cdots$\\
\bottomrule
\end{tabular}
\end{center}
\end{insight}

\paragraph{弱变分与强变分：邻域的两种度量。}
函数空间上可定义不同范数，导致不同强度的极值概念。\emph{弱变分}用 $C^1$ 范数 $\|h\|_1=\max|h|+\max|h'|$（函数值与导数都小，只比较「形状相近」的竞争者）；\emph{强变分}用 $C^0$ 范数 $\|h\|_0=\max|h|$（只要函数值小，导数可任意大，允许「锯齿形」竞争者）。这一区分并非吹毛求疵：最优控制中的 Bang--Bang 控制（\cref{sec:pmp-bangbang}）正是强变分邻域里的竞争者——它在所有「光滑邻近控制」中最优（弱极小），却要靠强变分才能与「急速切换」的竞争者比较。EL 方程是两者共同的必要条件，但充分条件不同：弱极小只需 Legendre 条件，强极小还需 Weierstrass 条件（\cref{sec:pmp-second-order}）。

\begin{definition}[Gateaux 导数 / 第一变分]\label{def:pmp-first-variation}
泛函 $J$ 在 $y$ 处沿方向 $h$ 的 Gateaux 导数（第一变分）为
\begin{equation}
  \delta J[y;h]=\lim_{\varepsilon\to 0}\frac{J[y+\varepsilon h]-J[y]}{\varepsilon}
  =\left.\frac{\mathrm d}{\mathrm d\varepsilon}\right|_{\varepsilon=0}J[y+\varepsilon h].
  \label{eq:pmp-gateaux}
\end{equation}
若存在连续线性泛函 $DJ_y$ 使 $J[y+h]=J[y]+DJ_y(h)+o(\|h\|)$，则称 $J$ 在 $y$ 处 Fréchet 可微。Fréchet 可微蕴含 Gateaux 可微，反之不然。
\end{definition}

变分法几乎只用第一变分（Gateaux 导数）：我们只需「对所有方向第一变分为零」这一条件，而它在 Fréchet 导数可能不存在的非线性问题中仍然成立——这正是变分法适用面更广的原因。极值的一阶必要条件由此立得：若 $y^\ast$ 是（弱或强）极小，令 $\Phi(\varepsilon)=J[y^\ast+\varepsilon h]$，则 $\Phi(0)\le\Phi(\varepsilon)$ 对小 $|\varepsilon|$ 成立，故 $\Phi'(0)=\delta J[y^\ast;h]=0$ 对所有容许 $h$。这与有限维 $\nabla f(\mathbf x^\ast)=\mathbf 0$ 完全平行（\cref{thm:oc-variation-zero} 已陈述此结论，这里给出其泛函分析背景）。

\section{Euler--Lagrange 方程：从积分条件到逐点 ODE}
\label{sec:pmp-el}

第一变分为零是关于「所有方向 $h$」的条件。如何从中提取 $y^\ast$ 本身必须满足的微分方程？逻辑结构是：先算出第一变分的显式表达式，再用变分基本引理把「积分为零」化为「被积函数为零」。

\begin{theorem}[Euler--Lagrange 方程]\label{thm:pmp-el}
设 $y^\ast\in C^2[a,b]$ 是 \cref{eq:pmp-integral-functional} 在固定端点 $y(a)=A,\,y(b)=B$ 下的极值，$L\in C^2$。则
\begin{equation}
  \frac{\partial L}{\partial y}-\frac{\mathrm d}{\mathrm dx}\frac{\partial L}{\partial y'}=0.
  \label{eq:pmp-el}
\end{equation}
\end{theorem}

\begin{proof}
\textbf{(i) 参数化扰动.}\;取 $\eta\in C^2[a,b]$ 满足 $\eta(a)=\eta(b)=0$（保证扰动后仍在容许类），令 $y_\varepsilon=y^\ast+\varepsilon\eta$。
\textbf{(ii) 第一变分.}\;由 Leibniz 法则（需 $L$ 连续可微以交换求导与积分）
\begin{equation}
  \delta J=\int_a^b\Big[L_y\,\eta+L_{y'}\,\eta'\Big]\mathrm dx,
  \label{eq:pmp-el-first-var}
\end{equation}
其中 $L_y,L_{y'}$ 在 $y^\ast$ 处取值。
\textbf{(iii) 分部积分.}\;对第二项,
$\int_a^b L_{y'}\eta'\,\mathrm dx=[L_{y'}\eta]_a^b-\int_a^b\big(\tfrac{\mathrm d}{\mathrm dx}L_{y'}\big)\eta\,\mathrm dx$；
由 $\eta(a)=\eta(b)=0$ 边界项消失。
\textbf{(iv) 合并.}\;$\delta J=\int_a^b\big[L_y-\tfrac{\mathrm d}{\mathrm dx}L_{y'}\big]\eta\,\mathrm dx=0$，对所有满足零端点的 $\eta$ 成立。
\textbf{(v) 变分引理}（见 \cref{lem:pmp-dubois}）将积分条件化为逐点条件，得 \cref{eq:pmp-el}。
\end{proof}

\begin{lemma}[变分基本引理 / du Bois-Reymond]\label{lem:pmp-dubois}
若 $f\in C[a,b]$，且对所有 $\eta\in C_0^\infty(a,b)$ 有 $\int_a^b f\eta\,\mathrm dx=0$，则 $f\equiv 0$。
\end{lemma}

\begin{proof}
反设 $f(x_0)>0$（$<0$ 类似）。由连续性，存在 $\delta>0$ 使 $f>0$ 于 $(x_0-\delta,x_0+\delta)\subset(a,b)$。取凸起（bump）函数
\begin{equation}
  \eta(x)=\begin{cases}\exp\!\big(-1/[\delta^2-(x-x_0)^2]\big),&|x-x_0|<\delta,\\[2pt]0,&\text{否则,}\end{cases}
  \label{eq:pmp-bump}
\end{equation}
则 $\eta\in C_0^\infty$、$\eta\ge 0$、在 $(x_0-\delta,x_0+\delta)$ 上严格正，于是 $\int f\eta\,\mathrm dx>0$，矛盾。本质上这是 $C_0^\infty$ 在 $L^2$ 中稠密——测试函数族足够丰富，可探测任一点的非零性。du Bois-Reymond（1879）进一步证明：仅设 $f$ 可积，结论仍以「几乎处处为零」成立，这是现代弱解理论的先驱。
\end{proof}

\paragraph{物理意义与展开。}
EL 方程可改写为 $\tfrac{\mathrm d}{\mathrm dx}L_{y'}=L_y$：左边是「广义动量的变化率」，右边是「广义力」，恰是 Newton 第二定律 $\dot p=F$ 的泛函版本。把 $\tfrac{\mathrm d}{\mathrm dx}L_{y'}$ 用链式法则展开
\begin{equation}
  L_y-L_{y'x}-L_{y'y}\,y'-L_{y'y'}\,y''=0,
  \label{eq:pmp-el-expanded}
\end{equation}
这是关于 $y$ 的二阶 ODE（设 $L_{y'y'}\ne 0$），配两个端点条件构成两点边值问题（BVP）。若 $L_{y'y'}=0$（$L$ 关于 $y'$ 线性），方程退化——这在最优控制中正对应\emph{奇异弧}（\cref{sec:pmp-bangbang}）。若 $y^\ast$ 只 $C^1$（有「角点」），EL 在角点失效，代之以 Weierstrass--Erdmann 角条件（\cref{sec:pmp-constraints}）。

\paragraph{多维系统与高阶情形。}
当 $\mathbf y=(y_1,\dots,y_n)$ 时得 $n$ 个耦合方程 $\tfrac{\partial L}{\partial y_i}-\tfrac{\mathrm d}{\mathrm dx}\tfrac{\partial L}{\partial y_i'}=0$。取 $x=t,\ \mathbf y=\mathbf q,\ L=T-V$ 即 Lagrange 力学（\cref{sec:pmp-hamilton}）。含高阶导数 $L(x,y,\dots,y^{(n)})$ 时反复分部积分得 \emph{Euler--Poisson 方程}
\begin{equation}
  \sum_{k=0}^{n}(-1)^k\frac{\mathrm d^k}{\mathrm dx^k}\frac{\partial L}{\partial y^{(k)}}=0,
  \label{eq:pmp-euler-poisson}
\end{equation}
符号交替 $(-1)^k$ 来自每次分部积分的负号。这直接决定了平滑轨迹的多项式阶数：minimum-jerk（$L=\dddot x^2$）给六阶 ODE、解为五次多项式；minimum-snap（$L=x^{(4)\,2}$）给八阶 ODE、解为七次多项式——后者是 Mellinger--Kumar 四旋翼微分平坦轨迹规划\cite{mellinger2011minimum} 的数学核心。

\begin{pitfall}[概念误区：EL 方程的解一定是极小]
EL 方程只是必要条件，类比有限维 $\nabla f=\mathbf 0$，其解可能是极小、极大或鞍点。求出 EL 解后\textbf{必须}用二阶条件（Legendre、Jacobi、Weierstrass，\cref{sec:pmp-second-order}）验证。另一常见错误是分部积分时忘记检查边界项：自由端点时边界项恰恰给出「自然边界条件」，是横截条件的来源。
\end{pitfall}

\section{经典变分问题：EL 与第一积分实战}
\label{sec:pmp-classics}

EL 方程是通用工具，但要真正掌握它须在具体问题上反复演练。一个反复奏效的简化是\emph{第一积分}——利用 Lagrangian 的对称性把二阶 ODE 降为一阶。

\begin{theorem}[Beltrami 恒等式]\label{thm:pmp-beltrami}
若 $L=L(y,y')$ 不显含 $x$，则沿 EL 解 $L-y'L_{y'}=C$（常数）。
\end{theorem}

\begin{proof}
$\tfrac{\mathrm d}{\mathrm dx}(L-y'L_{y'})=L_yy'+L_{y'}y''-y''L_{y'}-y'\tfrac{\mathrm d}{\mathrm dx}L_{y'}=y'\big(L_y-\tfrac{\mathrm d}{\mathrm dx}L_{y'}\big)=0$，
末步用 \cref{eq:pmp-el}。
\end{proof}

对力学系统 $L=T(\dot q)-V(q)$，Beltrami 给出 $L-\dot qL_{\dot q}=-(T+V)=-H$，即\emph{能量守恒}——这是 \cref{sec:pmp-hamilton} Noether 定理在时间平移对称下的特例，也预示了 PMP 的「Hamiltonian 沿最优轨迹守恒」（\cref{sec:pmp-statement}）。类似地，$L=L(x,y')$（不含 $y$）给出 $L_{y'}=\text{const}$，即「循环坐标 $\Rightarrow$ 动量守恒」。统一的工作流是
\begin{equation}
  \text{泛函极值}\;\xrightarrow{\delta J=0}\;\text{EL}\;\xrightarrow{\text{对称性/第一积分}}\;\text{降阶 ODE}\;\xrightarrow{\text{积分}}\;\text{解.}
  \label{eq:pmp-workflow}
\end{equation}

\paragraph{最速降线（摆线）。}
质点从 $A=(0,0)$ 受重力沿 $y(x)$（$y$ 轴向下为正）滑到 $B$，最小化时间。由 $v=\sqrt{2gy}$、$\mathrm ds=\sqrt{1+y'^2}\,\mathrm dx$，
$T[y]=\int_0^{x_1}\sqrt{(1+y'^2)/(2gy)}\,\mathrm dx$，$L$ 不含 $x$。Beltrami 给 $1/[\,2gy(1+y'^2)\,]=C^2$，即 $y(1+y'^2)=2r$。参数化 $y'=\cot(\theta/2)$ 得
\begin{equation}
  x=r(\theta-\sin\theta),\qquad y=r(1-\cos\theta),
  \label{eq:pmp-cycloid}
\end{equation}
即\emph{摆线}（圆在直线上滚动时圆周一点的轨迹）。物理直觉：曲线初段陡峭使质点迅速获得高速，「先加速再匀速」总时间短于「匀加速走直线」。这道 1696 年 Johann Bernoulli 的挑战题催生了整个变分法，六位巨匠给出解答（Johann/Jakob Bernoulli、Newton、Leibniz、L'Hôpital、Tschirnhaus）\cite{sussmann1997threehundred}。

\begin{figure}[ht]
\centering
\begin{tikzpicture}[>=Stealth, scale=1.0]
  % cycloid x=t-sin t, y=-(1-cos t), scaled
  \draw[->] (-0.2,0) -- (5.4,0) node[right]{$x$};
  \draw[->] (0,0.3) -- (0,-2.6) node[below]{$y$};
  \fill (0,0) circle (1.6pt) node[above left]{$A$};
  \draw[very thick, main] plot[domain=0:3.14159,samples=60,variable=\t]
       ({\t-sin(\t r)},{-(1-cos(\t r))});
  \fill ({3.14159},{-2}) circle (1.6pt) node[below right]{$B$};
  \draw[dashed, gray] (0,0) -- ({3.14159},{-2}) node[midway,above,sloped]{\small 直线（更慢）};
  \node[main] at (1.4,-1.55) {\small 摆线（最速降线）};
\end{tikzpicture}
\caption{最速降线是摆线：初段陡峭以快速积累速度，全局用时短于直线。}
\label{fig:pmp-cycloid}
\end{figure}

\paragraph{悬链线与最小旋转面：存在性的教训。}
均匀绳索在重力下平衡使势能 $\int y\sqrt{1+y'^2}\,\mathrm dx$ 最小，Beltrami 给 $y/\sqrt{1+y'^2}=c$，解为\emph{悬链线} $y=c\cosh\!\frac{x-x_0}{c}$（Galileo 曾误猜为抛物线；二者仅在小下垂时近似）。把同一曲线绕轴旋转得最小旋转面，面积泛函 $A=2\pi\int y\sqrt{1+y'^2}\,\mathrm dx$ 形式相同，故解仍是悬链线（catenoid）。但当两圆环相距过远——临界比 $d/(r_1+r_2)\approx 0.6627$ 被突破时——光滑悬链线解\emph{不存在}，最小面积「解」退化为两个圆盘（Goldschmidt 不连续解）。这是变分法的重要教训：\textbf{EL 有解 $\ne$ 变分问题有解}；完整理论还需存在性（Tonelli 直接法）、正则性与全局比较。

\begin{note}
关于临界比的常见混淆：等半径环（$r_1=r_2=R$）常引「分离距离 $d\approx 1.325R$」，换算为 $d/(2R)\approx 0.6627$，与上式自洽。网上另见的 $0.5277$ 是\emph{另一}归一化（等面积阈），与此处的 $d/(r_1+r_2)$ 不是同一量，不应据以「订正」。
\end{note}

\paragraph{测地线。}
Riemann 流形上两点最短曲线用能量泛函 $E[\gamma]=\tfrac12\int g_{ij}\dot x^i\dot x^j\,\mathrm dt$（避免长度泛函的参数化退化）。EL 经 Christoffel 符号 $\Gamma^k_{ij}=\tfrac12 g^{kl}(\partial_ig_{jl}+\partial_jg_{il}-\partial_lg_{ij})$ 化简得
\begin{equation}
  \ddot x^k+\Gamma^k_{ij}\dot x^i\dot x^j=0.
  \label{eq:pmp-geodesic}
\end{equation}
$\SO(3)$ 上的测地线给出旋转的最短路径，是球面线性插值（SLERP）与 $\SE(3)$ 位姿插值的理论基础（与 \cref{ch:lie} 的指数映射、右扰动约定一致，亦见 \cref{ch:rigid_body}）。更深刻地，取构型空间度量 $g_{ij}=M_{ij}(\mathbf q)$（惯性矩阵）时测地线就是无外力的自由运动——这是「惯性矩阵即构型空间度量」的几何含义，与 \cref{eq:pmp-robot-dyn} 的 Coriolis 结构互为表里。等时降线（tautochrone）则是摆线的另一推论：倒摆线上质点滑到最低点的时间 $T=\pi\sqrt{r/g}$ 与释放高度无关，Huygens（1673）据此设计摆线摆钟。

\section{约束变分与横截条件：通往共态的概念主梁}
\label{sec:pmp-constraints}

实际极值函数往往受额外约束。乘子法的精妙在于：引入额外参数把约束优化化为无约束优化。约束的「全局/逐点」结构决定了乘子是常数还是函数——\textbf{这条线索正是 PMP 共态的前身}，本节反复点明。

\paragraph{等周约束：常数乘子。}
在积分约束 $K[y]=\int_a^b G\,\mathrm dx=\ell$ 下极小化 $J$，引入\emph{常数}乘子 $\lambda$，对增广 Lagrangian $\tilde L=L+\lambda G$ 写 EL 方程，配约束定 $\lambda$。乘子是常数，因为 $\int G\,\mathrm dx=\ell$ 是\emph{一个全局约束}。经典实例 Dido 问题（固定周长围最大面积）：$\tilde L=-y+\lambda\sqrt{1+y'^2}$，积分得圆弧 $(x-c_1)^2+(y-c_2)^2=\lambda^2$，复现等周不等式 $4\pi A\le \ell^2$ 取等于圆。乘子有清晰含义 $\lambda=\mathrm dJ^\ast/\mathrm d\ell$——约束松弛单位增量时最优值的变化率（经济学的影子价格、最优控制共态的边际值，三者本质相同；有限维对应见 \cref{ch:nlopt} 的 KKT 与拉格朗日对偶）。

\paragraph{Holonomic 约束：函数乘子。}
逐点约束 $\phi(x,y_1,\dots,y_n)=0$ 下，乘子升级为\emph{函数} $\lambda(x)$：
\begin{equation}
  \frac{\partial L}{\partial y_i}-\frac{\mathrm d}{\mathrm dx}\frac{\partial L}{\partial y_i'}+\lambda(x)\frac{\partial\phi}{\partial y_i}=0,\quad i=1,\dots,n,
  \label{eq:pmp-holonomic}
\end{equation}
配约束 $\phi=0$ 共 $n+1$ 个方程。乘子是函数，因为约束在\emph{每个} $x$ 都要满足，$\lambda(x)$ 是「在 $x$ 处维持约束所需的力」。

\begin{insight}[贯穿全章的概念主梁]
等周约束的乘子是\emph{常数}（全局/积分约束）；holonomic 约束的乘子是\emph{函数} $\lambda(x)$（逐点约束）。\textbf{这正是 PMP 共态 $\boldsymbol\lambda(t)$ 的前身}——共态满足一个 ODE、是「逐时刻」约束的乘子，对应动力学约束 $\dot{\mathbf x}=\mathbf f(\mathbf x,\mathbf u)$ 这个在每个时刻都成立的等式。本章后续把这条「常数乘子 $\to$ 函数乘子 $\to$ 共态」的线索一路贯通到 \cref{sec:pmp-statement}。
\end{insight}

\paragraph{非完整约束：变分原理的微妙。}
当约束涉及速度（如轮式机器人无侧滑约束 $\dot x\sin\theta-\dot y\cos\theta=0$，不可积，Frobenius 条件不满足），d'Alembert 原理（虚位移满足约束）与 Hamilton 原理（用乘子强约束 $\dot q$）给出\emph{不同}方程：前者得正确运动方程 $\mathbf M\ddot{\mathbf q}+\mathbf C\dot{\mathbf q}+\mathbf g=\mathbf B\boldsymbol\tau+\mathbf A^\top\boldsymbol\lambda$（$\mathbf A\dot{\mathbf q}=0$），后者得一般与物理不符的「vakonomic」方程。这一陷阱在轮式、蛇形、滚动接触机器人中至关重要，约束动力学细节见 \cref{ch:constrained-dynamics}。

\paragraph{横截条件：自由端点的力平衡。}
端点不完全固定时需额外边界条件。\cref{ch:optimal-control-basics}（\cref{eq:oc-transversality-free}）已给基本情形，这里汇总三类并显式接到 PMP：
\begin{center}\small
\begin{tabular}{lll}
\toprule
变分法端点条件 & PMP 对应 & 共态终端条件\\
\midrule
$y(b)=B$（固定） & $\mathbf x(t_f)=\mathbf x_1$（全固定） & $\boldsymbol\lambda(t_f)$ 自由（由 BVP 定）\\
$y(b)$ 自由，无终端代价 & $\mathbf x(t_f)$ 自由 & $\boldsymbol\lambda(t_f)=\mathbf 0$\\
$y(b)$ 自由，终端代价 $\Phi$ & 有终端代价 $\phi$ & $\boldsymbol\lambda(t_f)=\nabla\Phi$\\
$y(b)\in\Gamma$ & $\mathbf x(t_f)\in M_T$ & $\boldsymbol\lambda(t_f)\perp T_{\mathbf x(t_f)}M_T$\\
\bottomrule
\end{tabular}
\end{center}
含终端代价时自然边界条件成 $L_{y'}|_b=-\,\mathrm d\Phi/\mathrm dy|_{y(b)}$，正是 PMP 横截条件 $\boldsymbol\lambda(t_f)=\nabla\Phi(\mathbf x(t_f))$ 的直接前身。横截条件本质是「边界处的力平衡」：极值函数在自由端不能有「剩余力」推它改变。

\paragraph{Weierstrass--Erdmann 角条件。}
若极值函数仅分段 $C^2$（在 $x=c$ 有角点），EL 在角点失效，代之以「动量连续」$L_{y'}(c,y,y'^-)=L_{y'}(c,y,y'^+)$ 与「Hamiltonian 连续」$[L-y'L_{y'}]^-=[L-y'L_{y'}]^+$。在光学中这是 Snell 折射定律（$n\sin\theta\leftrightarrow L_{y'}$ 连续）；在最优控制中对应 Bang--Bang 的\emph{切换条件}——切换瞬间共态与 Hamiltonian 必须连续。

\section{Hamilton 原理、Legendre 变换与辛结构}
\label{sec:pmp-hamilton}

变分法最深刻的应用是力学的 Lagrange/Hamilton 形式——它既是机器人动力学 $\mathbf M\ddot{\mathbf q}+\mathbf C\dot{\mathbf q}+\mathbf g=\boldsymbol\tau$ 的来源，又提供了 PMP 共态方程的原型。

\begin{theorem}[Hamilton 原理]\label{thm:pmp-hamilton}
$n$ 自由度系统取 $L=T(\mathbf q,\dot{\mathbf q})-V(\mathbf q)$，真实运动 $\mathbf q^\ast(t)$ 使作用量 $S[\mathbf q]=\int_{t_1}^{t_2}L\,\mathrm dt$ 取\emph{驻值}（$\delta S=0$），即满足 EL 方程。
\end{theorem}

\begin{pitfall}[「最小作用量」的误称]
Hamilton 原理说的是\emph{驻值}（stationary），不是\emph{最小}。短时间轨迹的作用量确取极小；但越过第一个共轭点（\cref{sec:pmp-second-order}）后可能只是鞍点。「最小作用量原理」沿用至今但数学上不准确——这也是 PMP「极值轨迹（extremal）$\ne$ 最优轨迹」警示的力学源头。
\end{pitfall}

对机器人取 $L=\tfrac12\dot{\mathbf q}^\top\mathbf M(\mathbf q)\dot{\mathbf q}-V(\mathbf q)$，EL 展开得
\begin{equation}
  \mathbf M(\mathbf q)\ddot{\mathbf q}+\mathbf C(\mathbf q,\dot{\mathbf q})\dot{\mathbf q}+\mathbf g(\mathbf q)=\boldsymbol\tau,\qquad
  C_{ij}=\sum_k\tfrac12\Big(\tfrac{\partial M_{ij}}{\partial q_k}+\tfrac{\partial M_{ik}}{\partial q_j}-\tfrac{\partial M_{jk}}{\partial q_i}\Big)\dot q_k,
  \label{eq:pmp-robot-dyn}
\end{equation}
$\mathbf C$ 的元素恰是惯性矩阵的 Christoffel 符号乘速度——印证「自由运动是构型空间测地线」。Lagrange 法相对 Newton 法的优势在约束系统中爆发：$n$ 连杆链用 Newton 须为每个连杆列力/力矩方程再消去约束力，方程数随 $n$ 爆炸；Lagrange 直接在广义坐标工作，约束力自动消失，方程数恰为自由度数。这正是 Pinocchio/RBDL 等库的理论根基；本章只给\emph{变分来源}，空间向量代数、RNEA/ABA 递推与几何力学细节见 \cref{ch:spatial-dynamics} 与 \cref{ch:geometric-mechanics}。

\paragraph{Legendre 变换与正则方程。}
最优控制更自然的变量是 $(\mathbf q,\mathbf p)$。定义共轭动量 $\mathbf p=\partial L/\partial\dot{\mathbf q}$、Hamiltonian $H(\mathbf q,\mathbf p,t)=\mathbf p^\top\dot{\mathbf q}-L$（$\dot{\mathbf q}$ 经 $\mathbf p=\partial L/\partial\dot{\mathbf q}$ 反解代入，可行性要求 Legendre 条件 $\det(\partial^2 L/\partial\dot q_i\partial\dot q_j)\ne0$）。EL 化为\emph{Hamilton 正则方程}
\begin{equation}
  \dot q_i=\frac{\partial H}{\partial p_i},\qquad \dot p_i=-\frac{\partial H}{\partial q_i},
  \label{eq:pmp-canonical}
\end{equation}
一个 $2n$ 维一阶系统。它的优美对称是辛几何核心，更是 PMP 共态方程 $\dot{\boldsymbol\lambda}=-\partial H/\partial\mathbf x$ 的直接原型——后者恰是正则方程的「一半」。对机器人 $\mathbf p=\mathbf M\dot{\mathbf q}$、$H=\tfrac12\mathbf p^\top\mathbf M^{-1}\mathbf p+V=T+V$（总能量）。

\paragraph{辛结构与 Noether 定理。}
记 $\mathbf z=(\mathbf q,\mathbf p)$ 与辛矩阵 $\boldsymbol\Omega=\begin{psmallmatrix}\mathbf 0&\mathbf I_n\\-\mathbf I_n&\mathbf 0\end{psmallmatrix}$，正则方程即 $\dot{\mathbf z}=\boldsymbol\Omega\nabla_{\mathbf z}H$。Hamilton 流保持辛形式 $\omega=\sum_i\mathrm dp_i\wedge\mathrm dq_i$，故相空间辛体积守恒（Liouville 定理）。这对长时间机器人仿真至关重要：普通 Runge--Kutta 不保辛、能量线性漂移，辛积分器（隐式中点、Störmer--Verlet）能量误差有界——辛积分器实现见 \cref{ch:geometric-mechanics}，本章仅点出 PMP 的「最优控制结构天然是辛的」这一观察。Poisson 括号 $\{F,G\}=\sum_i(\partial_{q_i}F\,\partial_{p_i}G-\partial_{p_i}F\,\partial_{q_i}G)$ 使 $\dot F=\{F,H\}$，两守恒量之括号仍守恒（Poisson 定理）。守恒量的系统来源是 Noether 定理：

\begin{theorem}[Noether 定理]\label{thm:pmp-noether}
设 $L(\mathbf q,\dot{\mathbf q},t)$ 在单参数连续变换 $q_i\mapsto q_i+\varepsilon\xi_i,\ t\mapsto t+\varepsilon\tau$ 下不变（至多差一个全导数），则沿 EL 解
\begin{equation}
  I=\sum_i\frac{\partial L}{\partial\dot q_i}\xi_i+\Big(L-\sum_i\frac{\partial L}{\partial\dot q_i}\dot q_i\Big)\tau=\text{const}.
  \label{eq:pmp-noether}
\end{equation}
\end{theorem}

\begin{proof}[证明骨架]
不变性意味 $\tfrac{\mathrm d}{\mathrm d\varepsilon}|_0 L(\mathbf q+\varepsilon\boldsymbol\xi,\dot{\mathbf q}+\varepsilon\dot{\boldsymbol\xi},t+\varepsilon\tau)=\tfrac{\mathrm dB}{\mathrm dt}$。展开左边并用 EL 方程 $L_{q_i}=\tfrac{\mathrm d}{\mathrm dt}L_{\dot q_i}$ 消项，可把全式写成某量的 $\tfrac{\mathrm d}{\mathrm dt}$，该量减去 $B$ 即守恒流 $I$。
\end{proof}

\begin{center}\small
\begin{tabular}{llll}
\toprule
对称性 & $(\xi,\tau)$ & 守恒量 & 物理含义\\
\midrule
时间平移 & $(0,1)$ & 能量 $H=T+V$ & 系统不显含时间\\
空间平移 & $(\delta_{ik},0)$ & 线动量 $p_k$ & 空间均匀\\
旋转不变 & $(\boldsymbol\omega\times\mathbf q,0)$ & 角动量 $\mathbf q\times\mathbf p$ & 空间各向同性\\
\bottomrule
\end{tabular}
\end{center}
例如二维中心力场 $L=\tfrac12 m(\dot x^2+\dot y^2)-V(r)$ 在旋转下不变，生成元 $\xi_x=-y,\xi_y=x$，代入 \cref{eq:pmp-noether}（$\tau=0$）得 $I=m(x\dot y-y\dot x)$——角动量守恒。Beltrami 恒等式正是 Noether 定理在时间平移对称下的特例（守恒量为能量）。Noether（1918）的工作\cite{noether1918invariante}（由 Klein 代为宣读，其课程曾以 Hilbert 名义开设）被誉为 20 世纪最重要的数学定理之一。

\section{二阶条件：何时真是极小}
\label{sec:pmp-second-order}

EL 只是必要条件。要确认 EL 解确为极小（而非极大或鞍点），须检验二阶条件——工程中尤为重要：不验充分条件，你可能正执行一条「最大代价」轨迹而不自知。

第二变分 $\delta^2 J=\int_a^b[Ph^2+2Qhh'+R(h')^2]\,\mathrm dx$，其中 $P=L_{yy},Q=L_{yy'},R=L_{y'y'}$（沿极值曲线）。极小的二阶必要条件是 $\delta^2 J\ge 0$ 对所有 $h(a)=h(b)=0$ 成立，类比有限维 Hessian 半正定。但无穷维「半正定」远难验证，需分层处理。

\begin{theorem}[Legendre 必要条件]\label{thm:pmp-legendre}
若 $y^\ast$ 是弱极小，则 $L_{y'y'}\ge 0$ 于 $[a,b]$。
\end{theorem}

\begin{proof}[证明思路]
取高频测试函数 $h_n=\tfrac1n\eta(x)\sin(n(x-x_0))$（$\eta$ 是 $x_0$ 处凸起函数）。则 $h_n\to 0$（$C^0,C^1$），而 $(h_n')^2\approx\eta^2\cos^2(\cdot)$。代入 $\delta^2 J$，由 Riemann--Lebesgue 引理 $\cos^2$ 均值趋 $\tfrac12$，得 $\lim_n\delta^2 J[y^\ast;h_n]=\tfrac12\int R\eta^2\mathrm dx$。若 $R(x_0)<0$ 则对小 $\delta$ 为负，矛盾。
\end{proof}

Legendre 是逐点条件，保证「高频方向」非负，但「低频方向」还需 Jacobi 条件——即使 $R\ge0$，若 $P$ 项足够负且 $h$ 变化缓慢，积分仍可能为负。沿极值曲线的\emph{Jacobi（accessory）方程}
\begin{equation}
  \frac{\mathrm d}{\mathrm dx}\big(L_{y'y'}\,u'\big)-\Big(L_{yy}-\frac{\mathrm d}{\mathrm dx}L_{yy'}\Big)u=0
  \label{eq:pmp-jacobi}
\end{equation}
的解 $u$（$u(a)=0,u'(a)\ne0$）在 $(a,b)$ 内首次为零的点称 $a$ 的\emph{共轭点}。

\begin{theorem}[Jacobi 充分条件]\label{thm:pmp-jacobi}
若 $L_{y'y'}>0$（强化 Legendre）且 $(a,b]$ 内无共轭点，则 $y^\ast$ 是严格弱极小。
\end{theorem}

共轭点的几何直觉：从 $a$ 出发的一族极值曲线初始发散，可能在某点重新汇聚（如透镜焦点），该点即共轭点；过共轭点后原曲线不再最优。经典例：从北极出发的经线（大圆弧）在南极汇聚，故长度超过半个大圆的测地线不再是最短路径——这正是 \cref{eq:pmp-geodesic} 测地线「半周后失最优」的二阶解释。

\paragraph{Weierstrass E-函数与强极小。}
Legendre--Jacobi 只保证弱极小（$C^1$ 邻域最优）。允许大导数变化（$C^0$ 邻域）时需 Weierstrass 超越函数
\begin{equation}
  \mathcal E(x,y,y',p)=L(x,y,p)-L(x,y,y')-(p-y')L_{y'}(x,y,y').
  \label{eq:pmp-weierstrass}
\end{equation}
$\mathcal E\ge0$ 等价于 $L$ 关于第三变量（斜率）凸——曲线在切线之上。对小偏差 $\mathcal E\approx\tfrac12 L_{y'y'}(p-y')^2$，故 Legendre（小偏差）是 Weierstrass（所有 $p$）的特例。若沿极值曲线 $\mathcal E\ge0$ 对所有 $p$ 成立且可嵌入斜率场，则 $y^\ast$ 是强极小。Bang--Bang 控制正需强极小条件——它的竞争者有大导数变化。这套二阶理论将在 \cref{sec:pmp-extensions} 以 Hamiltonian 形式（Kelley 条件、共轭点）再次出现，是判定 PMP 极值轨迹真伪的工具。

\begin{example}[一道兼有共轭点的例题]\label{ex:pmp-conjugate}
$J=\int_0^T[(y')^2-y^2]\,\mathrm dx$，$y(0)=y(T)=0$。EL 为 $y''+y=0$，通解 $y=A\sin x$；$y(T)=0$ 要求 $A\sin T=0$。Legendre：$L_{y'y'}=2>0$。Jacobi 方程沿 $y^\ast=0$ 为 $u''+u=0$，解 $u=\sin x$ 在 $x=\pi$ 首次为零，故 $\pi$ 是共轭点。结论：$T<\pi$ 时 $y^\ast=0$ 是严格弱极小；$T>\pi$ 时取 $h=\sin(\pi x/T)$ 算得 $\delta^2 J=T[(\pi/T)^2-1]<0$，$y^\ast=0$ \emph{不是}极小。
\end{example}

\section{从无约束到有约束：为什么需要 PMP}
\label{sec:pmp-why}

变分法的核心假设是\emph{控制无约束}——$\delta J=0$ 可通过 $\partial H/\partial\mathbf u=\mathbf 0$ 实现。但真实机器人控制永远有界：推力 $|\mathbf u|\le u_{\max}$、力矩限位、接触力非负 $\lambda_n\ge0$，甚至离散控制 $\mathbf u\in\{0,1\}$。当最优控制恰在约束边界（如满推力）时，$\partial H/\partial\mathbf u=\mathbf 0$ 无解——极值点不在内部！\cref{ch:optimal-control-basics} 已以此动机引入 Pontryagin 原理的陈述（\cref{thm:oc-pmp}），本章给出其严格证明与完整结构分类。

弱变分（EL 的 $\delta\mathbf u$ 小）在两种情形失效：其一，$\mathbf u^\ast$ 在 $U$ 边界，正向扰动出界、负向只给半空间信息；其二，$U$ 非凸或离散（如 $\{-1,+1\}$），根本无连续切方向。\emph{针状变分}反其道而行：不让扰动幅度小，而让它在\emph{时间}上小——在极短区间 $[\tau-\varepsilon,\tau]$ 把控制换成 $U$ 中\emph{任意}一点 $\mathbf v$（$\mathbf v-\mathbf u^\ast$ 可很大！），但作用时间 $\varepsilon$ 小，对轨迹影响仍是 $O(\varepsilon)$。这是 PMP 的天才所在。

\begin{insight}[直接法 vs 间接法：两种数值哲学]
求解最优控制有两条路。\emph{间接法}（「先优化再离散」）：先写 EL/PMP 的必要条件（ODE 与 BVP），再数值求解。\emph{直接法}（「先离散再优化」）：先把连续问题离散为有限维 NLP，送进求解器。变分法里 Ritz/Galerkin/配点法已体现这一分野（Ritz 在有限维子空间极小化泛函；Galerkin 要求 EL 残差与基函数正交；配点法要求 EL 在配点精确成立）。本章只讲其\emph{思想}：直接法的 KKT 乘子就是 PMP 共态的离散近似（\cref{sec:pmp-extensions}），故「直接法已隐式在解 PMP」。库实现（CasADi、acados、Crocoddyl）见 \cref{ch:mpc-advanced} 与 \cref{ch:ddp-ilqr}；采样类直接法的类比见 \cref{ch:mppi}。
\end{insight}

\section{Pontryagin 极小值原理：陈述与证明骨架}
\label{sec:pmp-statement}

\paragraph{标准形式与三形式等价。}
给定状态 $\mathbf x:[t_0,t_f]\to\mathbb R^n$、控制 $\mathbf u:[t_0,t_f]\to U$（$U$ 为 Lebesgue 可测紧集）、动力学 $\dot{\mathbf x}=\mathbf f(\mathbf x,\mathbf u,t),\ \mathbf x(t_0)=\mathbf x_0$、端点约束 $\boldsymbol\psi(\mathbf x(t_f),t_f)=\mathbf 0$。代价有三种经典形式：\emph{Bolza}（最一般）$J=\phi(\mathbf x(t_f),t_f)+\int_{t_0}^{t_f}L\,\mathrm dt$；\emph{Mayer}（仅末端）$J=\phi$；\emph{Lagrange}（仅积分）$J=\int L\,\mathrm dt$。三者经\emph{状态扩维}一阶等价：Lagrange $\to$ Mayer 只需引入 $\dot x_{n+1}=L,\ x_{n+1}(t_0)=0$，则 $J=x_{n+1}(t_f)$；Mayer 是 Bolza 取 $L\equiv0$ 的特例；Bolza $\to$ Mayer 同法。等价性保证「问题的数学描述不影响物理结论」，工程上 Bolza 最自然（同时表过程代价与末端精度）。

\paragraph{控制 Hamiltonian 与两约定。}
引入共态 $\boldsymbol\lambda(t)\in\mathbb R^n$（余切空间 $T^\ast_{\mathbf x}\mathbb R^n$ 的余向量）与异常乘子 $\lambda_0\ge0$，定义\emph{控制 Hamiltonian}
\begin{equation}
  H(\mathbf x,\mathbf u,\boldsymbol\lambda,\lambda_0,t)=\boldsymbol\lambda^\top\mathbf f(\mathbf x,\mathbf u,t)-\lambda_0 L(\mathbf x,\mathbf u,t).
  \label{eq:pmp-hamiltonian}
\end{equation}
$\boldsymbol\lambda^\top\mathbf f$ 度量「沿 $\mathbf f$ 运动对末端代价的改善率」，$-\lambda_0 L$ 是运行代价的即时贡献。选使 $H$ 最大的 $\mathbf u$，即在「改善末端」与「减少运行代价」间取最佳平衡。

\begin{note}[两种约定对照——读文献先看清]
\begin{center}\small
\begin{tabular}{lll}
\toprule
 & 极大化约定（Pontryagin 原始，本章） & 极小化约定（Kirk/Bertsekas，\cref{ch:control}）\\
\midrule
Hamiltonian & $H_{\max}=\boldsymbol\lambda^\top\mathbf f-\lambda_0 L$ & $H_{\min}=L+\boldsymbol\lambda^\top\mathbf f$\\
极值条件 & $\mathbf u^\ast=\arg\max_{\mathbf v\in U}H_{\max}$ & $\mathbf u^\ast=\arg\min_{\mathbf v\in U}H_{\min}$\\
共态方程 & $\dot{\boldsymbol\lambda}=-\partial H_{\max}/\partial\mathbf x$ & $\dot{\boldsymbol\lambda}=-\partial H_{\min}/\partial\mathbf x$\\
共态符号 & $\boldsymbol\lambda_{\max}$ & $\boldsymbol\lambda_{\min}=-\boldsymbol\lambda_{\max}=\nabla_{\mathbf x}V$\\
\bottomrule
\end{tabular}
\end{center}
两约定差一个负号，描述同一物理定律（如 $e^{\pm i\omega t}$）。\textbf{每个与 \cref{ch:control}/\cref{ch:dp-hjb}/\cref{ch:lqr-lqg-riccati} 的接口处换算 $\boldsymbol\lambda\to-\boldsymbol\lambda$。}\;\textbf{工程铁律}：无论哪种约定，最终反馈律恒为负反馈 $\mathbf u=-\mathbf K\mathbf x$（$\mathbf K=\mathbf R^{-1}\mathbf B^\top\mathbf P\succ0$）；若推出正反馈，必是约定混用。
\end{note}

\emph{正常极值}（$\lambda_0>0$，可归一 $\lambda_0=1$）：代价真正进入 $H$。\emph{异常极值}（$\lambda_0=0$）：代价完全消失，$H=\boldsymbol\lambda^\top\mathbf f$，极大化仅由动力学几何决定——此时无论代价是 $\int\!u^2$ 还是 $\int\!|u|$，极值轨迹都一样。非平凡性要求 $(\lambda_0,\boldsymbol\lambda)\ne(0,0)$，否则所有条件退化为 $0=0$。

\begin{theorem}[Pontryagin 极小值原理：五件套]\label{thm:pmp-main}
设 $(\mathbf x^\ast,\mathbf u^\ast)$ 是 Bolza 问题的局部最优解，$\mathbf f,L,\phi,\boldsymbol\psi$ 关于 $(\mathbf x,t)$ 连续可微，$U$ 任意（不要求凸）。则存在非平凡 $(\lambda_0,\boldsymbol\lambda(\cdot))$ 与末端乘子 $\boldsymbol\nu$ 使：

\noindent\textnormal{(1)} 状态方程 $\dot{\mathbf x}^\ast=\partial H/\partial\boldsymbol\lambda=\mathbf f$，$\mathbf x^\ast(t_0)=\mathbf x_0$；\\
\textnormal{(2)} 共态方程 $\displaystyle\dot{\boldsymbol\lambda}=-\frac{\partial H}{\partial\mathbf x}=-\Big(\frac{\partial\mathbf f}{\partial\mathbf x}\Big)^{\!\top}\boldsymbol\lambda+\lambda_0\Big(\frac{\partial L}{\partial\mathbf x}\Big)^{\!\top}$；\\
\textnormal{(3)} \textbf{极大化条件}（核心）：对几乎处处 $t$，
\begin{equation}
  \mathbf u^\ast(t)\in\arg\max_{\mathbf v\in U}H(\mathbf x^\ast(t),\mathbf v,\boldsymbol\lambda(t),\lambda_0,t);
  \label{eq:pmp-max}
\end{equation}
\textnormal{(4)} 横截条件 $\boldsymbol\lambda(t_f)=\lambda_0(\partial\phi/\partial\mathbf x)^\top+(\partial\boldsymbol\psi/\partial\mathbf x)^\top\boldsymbol\nu$；$t_f$ 自由时另有 $H(t_f)+\lambda_0\,\partial\phi/\partial t=0$；\\
\textnormal{(5)} 自治系统（$\mathbf f,L$ 不显含 $t$）时 $H$ 沿最优轨迹守恒，自由末端时间问题该常数为 $0$。
\end{theorem}

五件套合成一个两点边值问题：$\mathbf x$ 初值已知、$\boldsymbol\lambda$ 终值由 (4) 给出、$\mathbf x$ 终值由 $\boldsymbol\psi=0$ 定，通常无解析解。\textbf{PMP 是必要条件而非充分}——满足五件套的轨迹称\emph{Pontryagin 极值轨迹（extremal）}，未必最优（如 $f'(x_0)=0$ 仅说明驻点）；验证最优需 HJB 充分性、二阶条件或穷举比较（\cref{sec:pmp-reductions,sec:pmp-extensions}）。

\begin{figure}[ht]
\centering
\begin{tikzpicture}[>=Stealth, node distance=5mm and 9mm,
  blk/.style={draw, rounded corners, align=center, font=\small, inner sep=3pt, fill=main!6, draw=main!60!black},
  ln/.style={->, thick, gray!55!black}]
  \node[blk] (pmp) {PMP\\五件套};
  \node[blk, right=of pmp, yshift=12mm] (st) {(1) 状态方程\\$\mathbf x^\ast(t)$ 前向};
  \node[blk, right=of pmp, yshift=4mm]  (co) {(2) 共态方程\\$\boldsymbol\lambda(t)$ 反向};
  \node[blk, right=of pmp, yshift=-4mm] (mx) {(3) 极大化\\$\mathbf u^\ast(t)$ 逐时刻};
  \node[blk, right=of pmp, yshift=-12mm](tr) {(4) 横截\\$\boldsymbol\lambda(t_f)$};
  \node[blk, below=14mm of co] (hc) {(5) $H$ 守恒\\（验证/不变量）};
  \foreach \t in {st,co,mx,tr} {\draw[ln] (pmp) -- (\t);}
  \draw[ln] (pmp) |- (hc);
\end{tikzpicture}
\caption{PMP 五件套：状态前向、共态反向、控制逐时刻极大化、端点横截、Hamiltonian 守恒。三者合成 BVP。}
\label{fig:pmp-five}
\end{figure}

\paragraph{证明骨架：针状变分 $\to$ 可达凸锥 $\to$ 分离 $\to$ 共态 + 极大化。}
这是整个最优控制理论最重要单一结果的来由，分五步（研究生密度，技术细节略去）。

\emph{Step 1（针状扰动）.}\;对 $\mathbf u^\ast$ 的 Lebesgue 点 $\tau$ 与任意 $\mathbf v\in U$，令 $\mathbf u_\varepsilon$ 在 $[\tau-\varepsilon,\tau]$ 取 $\mathbf v$、余处取 $\mathbf u^\ast$。则 $\|\mathbf u_\varepsilon-\mathbf u^\ast\|_{L^\infty}$ 可大（强变分），但 $\|\cdot\|_{L^1}=O(\varepsilon)$。

\emph{Step 2（端点一阶增量）.}\;$\tau$ 时刻轨迹受冲击 $\mathbf x_\varepsilon(\tau)-\mathbf x^\ast(\tau)=\varepsilon\,\Delta\mathbf f+o(\varepsilon)$，$\Delta\mathbf f=\mathbf f(\mathbf x^\ast,\mathbf v,\tau)-\mathbf f(\mathbf x^\ast,\mathbf u^\ast,\tau)$。沿线性化方程 $\dot{\delta\mathbf x}=\mathbf A(t)\delta\mathbf x$（$\mathbf A=\partial\mathbf f/\partial\mathbf x$）传播到末端：$\delta\mathbf x(t_f)=\varepsilon\,\boldsymbol\Phi(t_f,\tau)\Delta\mathbf f+o(\varepsilon)$，$\boldsymbol\Phi$ 为状态转移矩阵（\cref{ch:statespace-basics}）。

\emph{Step 3（多针 $\to$ 可达凸锥）.}\;$k$ 个针的正组合给出
\begin{equation}
  K=\operatorname{cone}\big\{\boldsymbol\Phi(t_f,\tau)[\mathbf f(\mathbf x^\ast,\mathbf v,\tau)-\mathbf f(\mathbf x^\ast,\mathbf u^\ast,\tau)]:\tau\in(t_0,t_f),\,\mathbf v\in U\big\},
  \label{eq:pmp-cone}
\end{equation}
即可达集在 $\mathbf x^\ast(t_f)$ 的一阶凸近似（Boltyanskii 的「approximating cone / tent」）。

\emph{Step 4（凸锥分离）.}\;最优 $\Rightarrow$ 无可行扰动使代价严格下降 $\Rightarrow$ 可达锥 $K$ 与下降锥 $K_-=\{\delta\mathbf x:\nabla\phi\cdot\delta\mathbf x<0\}$ 内部不交。由 Minkowski 分离定理，存在非零 $(\lambda_0,\boldsymbol\lambda(t_f))$ 用一个超平面分开二者，$\boldsymbol\lambda(t_f)$ 即超平面法向。

\emph{Step 5（分离 $\to$ 共态 ODE + 极大化）.}\;沿时间反传 $\boldsymbol\lambda(t)=\boldsymbol\Phi(t_f,t)^\top\boldsymbol\lambda(t_f)$ 自动满足 $\dot{\boldsymbol\lambda}=-\mathbf A^\top\boldsymbol\lambda$。分离条件「$K$ 中每方向与 $\boldsymbol\lambda(t_f)$ 成非负角」化为 $\boldsymbol\lambda(\tau)^\top\mathbf f(\mathbf x^\ast,\mathbf u^\ast)\ge\boldsymbol\lambda(\tau)^\top\mathbf f(\mathbf x^\ast,\mathbf v)$，加上代价项即极大化条件 $H(\mathbf u^\ast)\ge H(\mathbf v)$。\hfill$\square$

\begin{insight}[针状变分为何不需 $U$ 凸]
PMP 的优美在于 $U$ \emph{不需}凸。凸锥 $K$ 是多针正组合\emph{自动}生成的——即使每个 $\mathbf v$ 只来自非凸 $U$ 的单点，正组合仍生成凸锥。弱变分需 $U$ 的切空间（要求光滑性），针状变分只需 $U$ 中\emph{存在}那个 $\mathbf v$。证明用到的工具：Minkowski 分离（分锥）、Gronwall 不等式（控制 $o(\varepsilon)$）、Lebesgue 微分定理（几乎处处是 Lebesgue 点）、状态转移矩阵（线性化传播）。其中 Gronwall 常数 $\exp(L_{\mathrm{Lip}}(t_f-t_0))$ 正是单次打靶病态性的根源——这把抽象证明与数值方法焊在一起。
\end{insight}

\begin{figure}[ht]
\centering
\begin{tikzpicture}[>=Stealth, scale=1.0]
  \coordinate (O) at (0,0);
  % reachable cone K (upper right)
  \fill[main!12] (O) -- (2.6,0.6) arc (13:62:2.7) -- cycle;
  \draw[main!70!black, thick] (O) -- (2.6,0.6);
  \draw[main!70!black, thick] (O) -- (1.27,2.4);
  \node[main!60!black] at (2.05,1.55) {$K$（可达锥）};
  % descent half-plane (left of separating line)
  \draw[second!70!black, thick] (-2.0,1.4) -- (2.0,-1.4);
  \node[second!70!black, align=center] at (-1.5,-0.6) {$K_-$\\下降方向};
  % normal lambda(tf)
  \draw[->, very thick, third!60!black] (O) -- (1.0,1.43) node[above right]{$\boldsymbol\lambda(t_f)$};
  \node[below left] at (O) {$\mathbf x^\ast(t_f)$};
\end{tikzpicture}
\caption{凸锥分离：可达锥 $K$ 与下降锥 $K_-$ 被超平面分开，法向即终端共态 $\boldsymbol\lambda(t_f)$，沿时间反传得共态 ODE。}
\label{fig:pmp-cone}
\end{figure}

\section{PMP 的三大归约：EL、HJB、LQR}
\label{sec:pmp-reductions}

PMP 不是孤立定理，它把变分法、动态规划、线性二次最优控制统一在一个框架下。本节给出三条归约——但严格遵守密度分工：HJB 主体留 \cref{ch:dp-hjb}、Riccati 理论留 \cref{ch:lqr-lqg-riccati}，这里只给 PMP 视角的\emph{接口}。

\paragraph{归约一：PMP $\Rightarrow$ EL。}
取 $\mathbf f=\mathbf u$、$U=\mathbb R^n$、$\lambda_0=1$，则 $H=\boldsymbol\lambda^\top\mathbf u-L$。无约束极大化 $\partial H/\partial\mathbf u=\boldsymbol\lambda-\partial L/\partial\mathbf u=0$ 给 $\boldsymbol\lambda=\partial L/\partial\dot{\mathbf x}$——\textbf{共态就是广义动量}。共态方程 $\dot{\boldsymbol\lambda}=-\partial H/\partial\mathbf x=\partial L/\partial\mathbf x$ 代入即 $\tfrac{\mathrm d}{\mathrm dt}\partial L/\partial\dot{\mathbf x}=\partial L/\partial\mathbf x$，正是 EL 方程；横截条件化为自然边界条件。故 PMP 严格推广 EL：EL 要求 $U=\mathbb R^m$ 且 $\mathbf f=\mathbf u$，PMP 只要 $U$ 紧、$\mathbf f$ 连续可微——这让所有欠驱动机器人（cart-pole、acrobot、四旋翼，$\dim\mathbf u\ll\dim\mathbf x$）落入 PMP 管辖。

\paragraph{归约二：PMP $\leftrightarrow$ HJB（必要 vs 充分）。}
值函数 $V(\mathbf x,t)=\inf_{\mathbf u(\cdot)}\{\int_t^{t_f}L\,\mathrm ds+\phi(\mathbf x(t_f))\}$ 满足 HJB 方程（极小化约定）
$\partial V/\partial t+\min_{\mathbf u\in U}\{L+\nabla_{\mathbf x}V^\top\mathbf f\}=0$，$V(\mathbf x,t_f)=\phi$。核心联系：沿 PMP 最优轨迹
\begin{equation}
  \boxed{\;\boldsymbol\lambda(t)=\nabla_{\mathbf x}V(\mathbf x^\ast(t),t)\;}\qquad(\text{极小化约定}),
  \label{eq:pmp-lambda-V}
\end{equation}
即 PMP 共态是 HJB 值函数沿最优轨迹的梯度，$H^\ast=-\partial V/\partial t$。证明骨架：对 HJB 关于 $\mathbf x$ 求导得 $\nabla_{\mathbf x}V$ 满足的 ODE，与共态方程同形且同终端条件 $\nabla_{\mathbf x}V(t_f)=\nabla\phi=\boldsymbol\lambda(t_f)$，由解唯一性得 \cref{eq:pmp-lambda-V}。

\begin{insight}[PMP 是 HJB 的特征线]
PMP 与 HJB 的关系精确等于一阶 PDE 与其特征线方程组的关系：HJB 是关于 $V(\mathbf x,t)$ 的一阶 PDE，PMP 的 $(\mathbf x(t),\boldsymbol\lambda(t))$ ODE 就是它的\emph{特征线}，沿特征线 PDE 降维为 ODE——这正是 PMP 不受维数灾难困扰的原因。代价是：PMP 只给必要条件，无法区分局部/全局最优；HJB 给充要条件但须求遍整个空间（维数灾难）。二者互补：PMP 给\emph{单条轨迹}的精确信息，HJB 给\emph{全状态空间}的最优反馈。\textbf{注意约定}：本章 PMP 用极大化、共态 $\boldsymbol\lambda_{\max}=-\nabla_{\mathbf x}V$；HJB 主体（黏性解、Crandall--Lions、特征线观点）见 \cref{ch:dp-hjb}，那里用极小化，接口处 $\boldsymbol\lambda\to-\boldsymbol\lambda$。
\end{insight}

\paragraph{归约三：PMP $\Rightarrow$ Riccati（LQR）。}
LQ 问题 $\min\tfrac12\mathbf x(t_f)^\top\mathbf Q_f\mathbf x(t_f)+\tfrac12\int(\mathbf x^\top\mathbf Q\mathbf x+\mathbf u^\top\mathbf R\mathbf u)\,\mathrm dt$，$\dot{\mathbf x}=\mathbf A\mathbf x+\mathbf B\mathbf u$，$\mathbf Q\succeq0,\mathbf R\succ0$。取 $\lambda_0=1$、极大化约定 $H=\boldsymbol\lambda^\top(\mathbf A\mathbf x+\mathbf B\mathbf u)-\tfrac12(\mathbf x^\top\mathbf Q\mathbf x+\mathbf u^\top\mathbf R\mathbf u)$。无约束极大化 $\partial H/\partial\mathbf u=\mathbf B^\top\boldsymbol\lambda-\mathbf R\mathbf u=0$ 给 $\mathbf u^\ast=\mathbf R^{-1}\mathbf B^\top\boldsymbol\lambda$（二阶 $\partial^2 H/\partial\mathbf u^2=-\mathbf R\prec0$ 确认极大）。代入得 Hamilton 矩阵系统
\begin{equation}
  \begin{pmatrix}\dot{\mathbf x}\\\dot{\boldsymbol\lambda}\end{pmatrix}=
  \underbrace{\begin{pmatrix}\mathbf A&\mathbf B\mathbf R^{-1}\mathbf B^\top\\ \mathbf Q&-\mathbf A^\top\end{pmatrix}}_{\mathcal H}
  \begin{pmatrix}\mathbf x\\\boldsymbol\lambda\end{pmatrix},
  \label{eq:pmp-hamilton-matrix}
\end{equation}
$\mathcal H$ 的特征值关于虚轴对称（辛结构）。\emph{Riccati 替换}令共态线性于状态，并直接取与本书工程约定一致的正定解记法 $\boldsymbol\lambda=-\tilde{\mathbf P}(t)\mathbf x$（极大化共态 $\boldsymbol\lambda_{\max}=-\nabla_{\mathbf x}V$，故 $\tilde{\mathbf P}\succeq0$）。对其求导并代入状态方程 $\dot{\mathbf x}=\mathbf A\mathbf x+\mathbf B\mathbf R^{-1}\mathbf B^\top\boldsymbol\lambda$ 与共态方程 $\dot{\boldsymbol\lambda}=\mathbf Q\mathbf x-\mathbf A^\top\boldsymbol\lambda$，整理（对任意 $\mathbf x$ 成立）得
\begin{equation}
  -\dot{\tilde{\mathbf P}}=\tilde{\mathbf P}\mathbf A+\mathbf A^\top\tilde{\mathbf P}-\tilde{\mathbf P}\mathbf B\mathbf R^{-1}\mathbf B^\top\tilde{\mathbf P}+\mathbf Q,\qquad \tilde{\mathbf P}(t_f)=\mathbf Q_f,
  \label{eq:pmp-riccati}
\end{equation}
即 Riccati 微分方程，终端条件来自横截条件 $\boldsymbol\lambda(t_f)=\mathbf Q_f\mathbf x(t_f)$（即 $\tilde{\mathbf P}(t_f)=\mathbf Q_f$）。最优控制 $\mathbf u^\ast=\mathbf R^{-1}\mathbf B^\top\boldsymbol\lambda=-\mathbf R^{-1}\mathbf B^\top\tilde{\mathbf P}\mathbf x=-\mathbf K\mathbf x$，恒为负反馈。（若改用记号 $\mathbf P:=\boldsymbol\lambda/\mathbf x=-\tilde{\mathbf P}$，则 \cref{eq:pmp-riccati} 的二次项与 $\mathbf Q$ 项同时反号，须留意勿与上式混淆。）

\begin{note}[换算到工程约定]
此处 $\tilde{\mathbf P}$ 即 \cref{ch:control}/\cref{ch:lqr-lqg-riccati} 的正定 Riccati 解，满足 $\boldsymbol\lambda_{\max}=-\nabla_{\mathbf x}V=-\tilde{\mathbf P}\mathbf x$；反馈律 $\mathbf u^\ast=-\mathbf R^{-1}\mathbf B^\top\tilde{\mathbf P}\mathbf x=-\mathbf K\mathbf x$ 为\textbf{负反馈}，与工程铁律一致。若有人按 $\mathbf P=\boldsymbol\lambda/\mathbf x$（$=-\tilde{\mathbf P}$，负定）写出 $\mathbf u^\ast=+\mathbf R^{-1}\mathbf B^\top\mathbf P\mathbf x$，那只是同一律的换号写法，绝非正反馈。本章只给「PMP $\Rightarrow$ Riccati 微分方程」这一步；DARE 存在唯一性、稳态 LQR、$(\mathbf A,\mathbf B)$ 可控 + $(\mathbf A,\mathbf Q^{1/2})$ 可观的渐近稳定性、分离原理见 \cref{ch:lqr-lqg-riccati}（亦与 \cref{ch:controllability-obsv,ch:linear-synthesis} 衔接）。LQR 之所以可解析，关键是动力学线性 + 代价二次使 $\boldsymbol\lambda=\mathbf P\mathbf x$ 保持线性，把无穷维 BVP 化为有限维 Riccati ODE。
\end{note}

\section{极大化结构与最优控制形态}
\label{sec:pmp-bangbang}

Hamiltonian 关于 $\mathbf u$ 的结构决定最优控制的形态。\cref{ch:optimal-control-basics}（\cref{thm:oc-bangbang}）已给最短时间 Bang--Bang 的相平面图解，本节做完整结构分类与切换理论。

\begin{center}\small
\begin{tabular}{lll}
\toprule
$H$ 关于 $\mathbf u$ & 极大化结果 & 典型问题\\
\midrule
线性 $H=a+\mathbf b^\top\mathbf u$ & Bang--Bang：$\mathbf u^\ast\in\partial U$ & 最短时间、推力受限\\
严格凹 $\partial^2 H/\partial\mathbf u^2\prec0$ & 内点解 $\partial H/\partial\mathbf u=0$ & LQR（$\mathbf u^\top\mathbf R\mathbf u$ 项）\\
退化 $\partial H/\partial\mathbf u\equiv0$ 区间 & 奇异弧（需高阶条件） & Goddard 中段\\
\bottomrule
\end{tabular}
\end{center}

\paragraph{Bang--Bang 与切换函数。}
当 $H=H_0(\mathbf x,\boldsymbol\lambda,t)+\boldsymbol\sigma^\top\mathbf u$（关于 $\mathbf u$ 仿射），极大化给
\begin{equation}
  u^\ast_i=\begin{cases}u_{i,\max},&\sigma_i>0,\\ u_{i,\min},&\sigma_i<0,\\ \text{待定},&\sigma_i=0,\end{cases}\qquad
  \boldsymbol\sigma=\frac{\partial H}{\partial\mathbf u}\ \text{（切换函数）}.
  \label{eq:pmp-switching}
\end{equation}
对线性时不变系统 $\dot{\mathbf x}=\mathbf A\mathbf x+\mathbf B\mathbf u$，$\boldsymbol\lambda$ 满足 $\dot{\boldsymbol\lambda}=-\mathbf A^\top\boldsymbol\lambda$，故 $\boldsymbol\sigma(t)=\mathbf B^\top e^{-\mathbf A^\top t}\boldsymbol\lambda(0)$ 是「多项式 $\times$ 指数」的组合，由 Chebyshev 系统理论其零点至多 $n-1$ 个：

\begin{theorem}[Feldbaum 切换上界]\label{thm:pmp-feldbaum}
对线性时不变系统 $\dot{\mathbf x}=\mathbf A\mathbf x+\mathbf B\mathbf u$、$\mathbf u\in[-1,1]^m$，最短时间控制的切换次数不超过 $n-1$\cite{feldbaum1953optimal}。
\end{theorem}

\begin{pitfall}[Bang--Bang 并非「粗暴」]
Bang--Bang 是理论上的\emph{精确}最优（对 $H$ 关于 $\mathbf u$ 线性的问题），不是近似或简化。航天器姿态机动、最短时间机械臂运动中，它比任何「平滑」控制更省时省燃料。一个示意的符号判定（伪代码，仅 $\le 3$ 行）：

\begin{center}\small\ttfamily
sigma = B' * lambda;\qquad u = sign(sigma);\qquad \%\ 切换发生在 sigma 过零处
\end{center}
\end{pitfall}

\paragraph{Fuller 问题：无穷切换反例。}
Feldbaum 的「最短时间」条件是本质的。对其他代价，切换可无穷。Fuller（1960）发现二维系统 $\dot x_1=x_2,\ \dot x_2=u,\ |u|\le1,\ J=\int_0^\infty x_1^2\,\mathrm dt$ 的最优控制有\emph{无穷次切换}（chattering）\cite{fuller1960relay}。其切换\emph{曲线}为 $x_1=-\operatorname{sign}(x_2)\,\xi\,x_2^2$，系数
\begin{equation}
  \xi=\sqrt{\tfrac{\sqrt{33}-1}{24}}\approx 0.4462,
  \label{eq:pmp-fuller-xi}
\end{equation}
而切换\emph{时刻}按几何级数收敛到零，相邻切换间隔之比为 $1/\alpha$，$\alpha=\sqrt{(1+2\xi)/(1-2\xi)}\approx 4.1301$，故比值 $\approx 0.242$。

\begin{note}[切换曲线系数 vs 时间公比，勿混]
常见误述把 $0.44$ 当作「切换时刻的几何公比」。\textbf{$0.4462$ 是切换曲线系数 $\xi$}（出现在切换曲线 $x_1=-\operatorname{sign}(x_2)\xi x_2^2$ 中）；\textbf{切换时刻的几何公比是 $1/\alpha\approx 1/4.13\approx 0.242$}。二者是不同的量，不可混用——这是「实测编造数值」类错误的典型。
\end{note}

\paragraph{奇异弧与 Kelley 条件。}
当切换函数 $\boldsymbol\sigma(t)\equiv0$ 于一整个区间（而非离散点），极大化退化——$H$ 对 $\mathbf u$ 是「平坦」的，一阶条件失效。这段称\emph{奇异弧}。确定奇异控制需逐阶求导直到 $\mathbf u$ 显式出现：$\sigma=0,\dot\sigma=0,\dots,\mathrm d^{2q}\sigma/\mathrm dt^{2q}=0$，最小整数 $q$ 称奇异弧的\emph{阶}，从末式解出 $u_{\mathrm{sing}}$。

\begin{theorem}[Kelley 广义 Legendre--Clebsch 条件]\label{thm:pmp-kelley}
标量控制奇异弧的 $q$ 阶必要条件为（本章极大化约定）
\begin{equation}
  (-1)^q\,\frac{\partial}{\partial u}\Big[\frac{\mathrm d^{2q}}{\mathrm dt^{2q}}\frac{\partial H}{\partial u}\Big]\le0.
  \label{eq:pmp-kelley}
\end{equation}
这是经典 Legendre 条件（$\partial^2 H/\partial u^2\le0$ 对极大化）的高阶推广——奇异弧上 $\partial^2 H/\partial u^2=0$ 退化，需看更高阶项\cite{kelley1964second}（向量推广见 Robbins\cite{robbins1967generalized}）。\footnote{符号随约定翻转：教材多以\emph{极小化}约定写作 $(-1)^q\,\partial_u[\mathrm d^{2q}\!/\mathrm dt^{2q}\,\partial_uH]\ge0$（$q=0$ 退回 $\partial^2H/\partial u^2\ge0$，对极小化）。本章用 Pontryagin 极大化约定，$q=0$ 须退回 $\partial^2H/\partial u^2\le0$，故整链不等号取 $\le0$（与 \cref{ch:control}/\cref{ch:lqr-lqg-riccati} 接口处随 $\boldsymbol\lambda\to-\boldsymbol\lambda$ 一并翻转）。极小化标准式见 Kelley(1964)\cite{kelley1964second}、Zelikin--Borisov 与 Bryson--Ho\cite{bryson1975applied}。}
\end{theorem}

奇异弧对直接法极不友好：离散化后 $\partial H/\partial u\approx0$ 让 NLP 的 Hessian 病态，典型现象是 Bang--Bang chattering（数值解在 $u_{\min},u_{\max}$ 间高频振荡「模拟」奇异弧）。间接法可行但须事先识别奇异结构、用同伦延拓跟踪。最省燃料问题 $J=\int|u|\,\mathrm dt$ 则给出第四种形态\emph{Bang-off-Bang}：$H=\lambda_1 x_2+\lambda_2 u-|u|$，极大化得 $u^\ast=+1\ (\lambda_2>1)$、$0\ (|\lambda_2|\le1)$、$-1\ (\lambda_2<-1)$——多出「滑行段 $u=0$」，体现「性能 vs 资源」的根本权衡。

\section{经典 PMP 实例：三种结构的范例}
\label{sec:pmp-examples}

\paragraph{实例一：双积分器最短时间（Bang--Bang）。}
$\ddot x=u,\ |u|\le1$，状态 $(x_1,x_2)=(x,\dot x)$，最小化到原点的 $t_f$。$L\equiv1$、$H=\lambda_1 x_2+\lambda_2 u-1$。极大化（$H$ 对 $u$ 线性）给 $u^\ast=\operatorname{sign}(\lambda_2)$。共态 $\dot\lambda_1=0,\ \dot\lambda_2=-\lambda_1$，故 $\lambda_2(t)$ 是一次多项式，至多一个零点——切换至多一次（Feldbaum，$n=2$）。$u=\pm1$ 下相轨迹为抛物线族 $x_1=\pm\tfrac12 x_2^2+c$，切换曲线
\begin{equation}
  \gamma:\ x_1=-\tfrac12 x_2|x_2|,\qquad
  u^\ast(x_1,x_2)=-\operatorname{sign}\!\big(x_1+\tfrac12 x_2|x_2|\big).
  \label{eq:pmp-switch-curve}
\end{equation}
曲线上方制动（$u=-1$）、下方加速（$u=+1$）；任意初态先沿一条抛物线到切换曲线，再沿曲线到原点（与 \cref{ch:optimal-control-basics} 的 \cref{eq:oc-switch-curve} 一致）。

\begin{figure}[ht]
\centering
\begin{tikzpicture}[>=Stealth, scale=0.95]
  \draw[->] (-3,0) -- (3,0) node[right]{$x_1$};
  \draw[->] (0,-2.2) -- (0,2.2) node[above]{$x_2$};
  % switching curve x1 = -1/2 x2|x2|
  \draw[very thick, second!70!black] plot[domain=0:2,samples=40,variable=\s] ({-0.5*\s*\s},{\s});
  \draw[very thick, second!70!black] plot[domain=0:2,samples=40,variable=\s] ({0.5*\s*\s},{-\s});
  \node[second!70!black] at (-1.7,1.65) {\small $\gamma$};
  % a sample trajectory: brake then slide
  \draw[->, thick, main] plot[domain=0:1.2,samples=30,variable=\t] ({2+\t-0.5*\t*\t},{1-\t});
  \draw[->, thick, main] plot[domain=0:1.34,samples=30,variable=\s] ({0.5*\s*\s},{-\s});
  \fill (2,1) circle (1.4pt) node[above right]{\small 初态};
  \fill (0,0) circle (1.4pt) node[below left]{\small $O$};
  \node[main] at (1.9,-0.2) {\small $u=-1$};
  \node[main] at (0.95,-1.55) {\small $u=+1$};
\end{tikzpicture}
\caption{双积分器最短时间的相平面综合：先沿一条抛物线（制动）到切换曲线 $\gamma$，再沿 $\gamma$（加速）到原点。}
\label{fig:pmp-bangbang}
\end{figure}

\paragraph{实例二：矩阵 LQR。}
已在 \cref{sec:pmp-reductions} 完整推出 $\mathbf u^\ast=-\mathbf R^{-1}\mathbf B^\top\tilde{\mathbf P}\mathbf x$，$\tilde{\mathbf P}$ 由 Riccati 反向积分。它是 PMP \emph{唯一}能手工解析解出的一般大类——动力学线性 + 代价二次同时满足，使共态与状态保持线性关系。所有非线性轨迹优化器（iLQR/DDP）的内核都是「沿名义轨迹局部 LQ 化 + 反向 Riccati 递推」（\cref{ch:ddp-ilqr}）。

\paragraph{实例三：Dubins 车六型。}
$\dot x=\cos\theta,\ \dot y=\sin\theta,\ \dot\theta=u,\ |u|\le1$。$H=\lambda_x\cos\theta+\lambda_y\sin\theta+\lambda_\theta u-1$。极大化给 $u^\ast=\operatorname{sign}(\lambda_\theta)$（Bang--Bang）。共态 $\dot\lambda_x=\dot\lambda_y=0$（守恒），$\dot\lambda_\theta=\lambda_x\sin\theta-\lambda_y\cos\theta$。弧段 $u=\pm1$ 上 $\theta$ 线性变化得圆弧（$L$/$R$），奇异弧 $\lambda_\theta\equiv0$ 给 $\theta=\text{const}$ 即直线段 $S$。共态零点分析表明最优路径至多三段，穷举得六型：

\begin{theorem}[Dubins 定理]\label{thm:pmp-dubins}
曲率受限的最短路径必为 $\{LSL,LSR,RSL,RSR,LRL,RLR\}$ 之一（$L$/$R$ 为最小半径左/右转弧、$S$ 为直线）\cite{dubins1957curves}。
\end{theorem}

这展示 PMP 最强大的工程价值之一：把「在所有曲线中找最短」的无穷维优化，降维为「在六型中选最好」的有限参数优化。Dubins 的证明本质就是 PMP 在该系统上的应用——共态结构（$\lambda_x,\lambda_y$ 守恒、$\lambda_\theta$ 是正弦函数）决定了路径类型的有限性。

\begin{note}[Reeds--Shepp：路径数的精确归属]
允许倒车（$v\in\{-1,+1\}$）时路径类型增多。Reeds--Shepp（1990）原文给出\textbf{48} 条充分路径集（48 words，归 CSC/CCC/CCCC/CCSC/CCSCC 五类）\cite{reeds1990optimal}；把 48 缩减为 \textbf{46} 的是 Sussmann--Tang（1991，用 Lie 代数方法）\cite{sussmann1991shortest}。常见误述把「46」错挂在 Reeds--Shepp 名下——正确表述是「Reeds--Shepp 1990 给 48、Sussmann--Tang 1991 缩到 46」。倒车入库的最短路径计算直接用 Reeds--Shepp 公式（固定翼无人机航路则用 Dubins，因不能倒飞）。
\end{note}

\paragraph{实例四：Goddard 火箭三段式。}
一维垂直 $\dot h=v,\ \dot v=(T-D)/m-g,\ \dot m=-T/c,\ 0\le T\le T_{\max}$，最大化 $h(t_f)$（Mayer）。$H$ 关于 $T$ 线性，切换函数 $\Sigma=\lambda_v/m-\lambda_m/c$。最优推力呈\emph{三段式}：$T_{\max}\to$ 奇异弧 $\to0$（coast）——$\Sigma>0$ 全推力、$\Sigma=0$ 奇异（$T_{\mathrm{sing}}$ 由 $\ddot\Sigma=0$ 定）、$\Sigma<0$ 滑行\cite{bonnans2008singular}。这是奇异弧最经典的工程实例。

\paragraph{实例五：四旋翼最短时间。}
现代机器人应用：Foehn--Romero--Scaramuzza（2021）用互补进度约束（CPC）把航点通过时刻作决策变量，数值求解 PMP 极值条件\cite{foehn2021cpc}。关键发现与 PMP 预测吻合：推力在绝大多数时段\emph{饱和}（Bang--Bang），航点附近现准奇异弧（协调三轴力矩的过渡段），实测击败职业竞速飞手。这印证了 \cref{sec:pmp-bangbang} 的结构分类在真实高维系统中的有效性。

\section{PMP 的现代扩展：理论选讲}
\label{sec:pmp-extensions}

本节是控制深度的理论选讲，把 PMP 推向状态约束、离散时间、混合系统、李群与充分性——这些是连接 \cref{ch:hj-reachability,ch:ddp-ilqr,ch:mpc-advanced} 与机械臂/动力学部的钩子。

\paragraph{状态约束 PMP：共态跳跃。}
实际系统除控制约束外常有状态约束 $g(\mathbf x,t)\le0$（避障 $\|\mathbf x-\mathbf o\|\ge r$、关节限位、地面不穿透）。此时共态 $\boldsymbol\lambda$ 可\emph{不连续}（跳跃）。引入 Borel 测度乘子 $\mu\in\mathcal M([t_0,t_f])$，$\operatorname{supp}(\mu)\subset\{t:g=0\}$、$\mu\ge0$（互补松弛），修正共态方程（分布意义）
\begin{equation}
  \mathrm d\boldsymbol\lambda=-\frac{\partial H}{\partial\mathbf x}\,\mathrm dt+\Big(\frac{\partial g}{\partial\mathbf x}\Big)^{\!\top}\mathrm d\mu.
  \label{eq:pmp-state-constraint}
\end{equation}
边界弧（沿 $g=0$ 滑行，$\mu$ 有连续密度）与接触点（瞬时触碰，$\mu$ 是 Dirac 质量、$\boldsymbol\lambda$ 在该点跳跃）需分别处理\cite{hartl1995survey}。约束\emph{阶}（$g$ 对时间求导直到 $\mathbf u$ 显现的次数）越高，边界弧控制越间接、数值越难。这与 CLF/CBF 安全滤波（\cref{ch:clf-cbf}）的状态约束概念相邻、与可达集（\cref{ch:hj-reachability}）的状态约束 $\to$ 可达集互引；接触约束的摩擦锥/LCP 见 \cref{ch:contact-mechanics}。

\paragraph{离散 PMP $\leftrightarrow$ DDP backward pass。}
连续 PMP 须离散实现。对 $\mathbf x_{k+1}=\mathbf f_k(\mathbf x_k,\mathbf u_k),\ J=\phi(\mathbf x_N)+\sum L_k$，存在共态序列使
\begin{equation}
  \boldsymbol\lambda_k=\Big(\frac{\partial\mathbf f_k}{\partial\mathbf x_k}\Big)^{\!\top}\boldsymbol\lambda_{k+1}+\Big(\frac{\partial L_k}{\partial\mathbf x_k}\Big)^{\!\top},\quad
  \boldsymbol\lambda_N=\Big(\frac{\partial\phi}{\partial\mathbf x_N}\Big)^{\!\top},\quad
  \mathbf u_k^\ast\in\arg\max_{\mathbf u\in U_k}\big\{\boldsymbol\lambda_{k+1}^\top\mathbf f_k-L_k\big\}
  \label{eq:pmp-discrete}
\end{equation}
（Halkin 1966\cite{halkin1966maximum}）。这精确对应 DDP backward pass：$Q$-函数展开 $Q_x=\ell_x+\mathbf f_x^\top V'_x,\ Q_u=\ell_u+\mathbf f_u^\top V'_x$ 中 $V'_x$ 正是下一时步共态 $\boldsymbol\lambda_{k+1}$（Jacobson--Mayne 1970\cite{jacobson1970differential}）；iLQR 丢弃 $V'_x\cdot\mathbf f_{xx}$ 项即 Gauss--Newton 近似。

\begin{pitfall}[共态滞后一步：$\boldsymbol\lambda_{k+1}$ 而非 $\boldsymbol\lambda_k$]
\cref{eq:pmp-discrete} 中出现在 $\mathbf u_k$ 极大化里的是 $\boldsymbol\lambda_{k+1}$（\emph{不是} $\boldsymbol\lambda_k$）——共态在时间上「滞后一步」。错误代码 \texttt{lambda[k]=fx'*lambda[k]+lx} 不会编译报错，却让轨迹优化完全不收敛或收敛到错误结果；正确为 \texttt{lambda[k]=fx'*lambda[k+1]+lx}。DDP 代码中这对应用 \texttt{Vx\_next}（下一时步值函数梯度）算当前控制。DDP/iLQR 主体与约束处理（ALTRO/Crocoddyl 谱系，库实现按本部范围剔除）见 \cref{ch:ddp-ilqr}。
\end{pitfall}

\paragraph{PMP $\leftrightarrow$ KKT：直接法已隐式解 PMP。}
离散化的最优控制是以 $(\mathbf x_k,\mathbf u_k)$ 为变量的 NLP。对动力学约束引入乘子 $\boldsymbol\mu_k$，KKT 条件对 $\mathbf x_k$ 求导后令 $\boldsymbol\lambda_k:=\boldsymbol\mu_{k-1}/\Delta t$，恰得 \cref{eq:pmp-discrete} 的离散共态方程——\textbf{NLP 乘子就是 PMP 共态的离散版本}，KKT 互补松弛就是控制约束的极大化条件。故用 CasADi+IPOPT 求轨迹时，即便从不显式写 PMP，求解器已在隐式解它；理解 PMP 的价值在于：你知道解该长什么样（Bang--Bang？奇异弧？切换几次？），从而判断数值解是否合理。这是实时迭代 MPC（\cref{ch:mpc-advanced}）的理论根，有限维 KKT/LICQ 见 \cref{ch:nlopt}。

\paragraph{数值间接法的病态性。}
间接法（先优化再离散）写出 $(\mathbf x,\boldsymbol\lambda)$ 的 BVP，猜 $\boldsymbol\lambda(t_0)$ 用牛顿法打靶。\emph{单次打靶}极度病态：$\boldsymbol\lambda(t_0)$ 的 $10^{-6}$ 扰动可在 $t_f$ 处放大为巨大误差——根源是 \cref{sec:pmp-statement} 的 Gronwall 常数 $\exp(L_{\mathrm{Lip}}(t_f-t_0))$，线性化系统的正 Lyapunov 指数使前向积分指数放大。对策是\emph{多重打靶}（分段、段间匹配条件作约束）或\emph{配点法}（所有网格点同时满足配点方程），并用同伦延拓提供初值。这里只给「方法存在性 + 病态性原因」的理论说明，调参经验与软件生态属工程实现。

\paragraph{混合系统 PMP：腿足步态。}
连续动力学 + 离散模式切换 $\dot{\mathbf x}=\mathbf f_{q(t)}(\mathbf x,\mathbf u)$ 在腿足机器人中是核心结构（每条腿的接触/脱离是离散模式）。在切换时刻 $t_s$：Hamiltonian 连续 $H_{q^-}(t_s^-)=H_{q^+}(t_s^+)$，共态跳跃 $\boldsymbol\lambda(t_s^-)=(\partial\Phi/\partial\mathbf x)^\top\boldsymbol\lambda(t_s^+)+p\,\nabla S$（$\Phi$ 为 reset map、$S$ 为切换面、$p$ 为切换面乘子）\cite{sussmann1999maximum,shaikh2007hybrid}。支撑相 $\lambda_n>0$（法向力存在）、摆动相 $\lambda_n=0$，切换在着地/离地瞬间。这是 \cref{eq:pmp-state-constraint} 状态约束与角条件（\cref{sec:pmp-constraints}）的混合系统版本，接触动力学见 \cref{ch:contact-mechanics}、全身控制见 \cref{ch:force-wbc}。

\paragraph{$\SE(3)$/Lie--Poisson 上的 PMP。}
四旋翼姿态、航天器机动、机械臂末端位姿的状态空间是李群而非 $\mathbb R^n$，标准 PMP 失效（无全局坐标、$\mathbf x+\delta\mathbf x$ 无定义）。坐标无关 PMP（Jurdjevic；Agrachev--Sachkov）将状态置于 $G$（左不变向量场 $\dot g=g\cdot\boldsymbol\xi$）、共态置于 $\mathfrak g^\ast$，服从 Lie--Poisson 方程 $\dot{\mathbf p}=\operatorname{ad}^\ast_{\boldsymbol\xi}\mathbf p+(\text{控制力矩对偶})$\cite{saccon2013optimal}。对自由刚体 $\mathbf p=\mathbf I\boldsymbol\omega$（角动量）满足 $\dot{\mathbf p}=\mathbf p\times\boldsymbol\omega$——正是 Euler 刚体方程，其 Lie--Poisson 结构守恒角动量模长。这里 $\operatorname{ad}^\ast$ 是 \cref{ch:lie} 伴随表示 $\operatorname{ad}$ 的对偶，须对齐本书\emph{右扰动}约定与指数/对数映射；离散几何最优控制（保辛李群积分器）见 Kobilarov--Marsden\cite{kobilarov2011discrete} 与 \cref{ch:geometric-mechanics}。

\paragraph{充分性：Mangasarian 与 LQR 全局最优。}
PMP 何时给全局最优？

\begin{theorem}[Mangasarian 充分条件]\label{thm:pmp-mangasarian}
若 $H(\mathbf x,\mathbf u,\boldsymbol\lambda)$ 关于 $(\mathbf x,\mathbf u)$ 联合凹（固定 $\boldsymbol\lambda$）、终端代价 $\phi$ 关于 $\mathbf x$ 凸，则满足 PMP 的极值轨迹是\emph{全局最优}\cite{mangasarian1966sufficient}。
\end{theorem}

LQR 恰满足：$-\tfrac12(\mathbf x^\top\mathbf Q\mathbf x+\mathbf u^\top\mathbf R\mathbf u)$ 凹、$\boldsymbol\lambda^\top(\mathbf A\mathbf x+\mathbf B\mathbf u)$ 线性，故 LQR 的 Riccati 反馈\emph{自动}全局最优、无需额外验证——这正是 \cref{sec:pmp-reductions} LQR 例的「全局性」来由。二阶充分条件（共轭点理论）的 Hamiltonian 版本：极值轨迹在 $[t_0,t_f]$ 无共轭点且强化 Legendre--Clebsch $\partial^2 H/\partial\mathbf u^2\prec0$ 成立，则局部严格最优；共轭点对应 Jacobi 方程的矩阵 Riccati 解发散，是 \cref{sec:pmp-second-order} 共轭点理论的最优控制对应。非光滑动力学（接触/摩擦）则需 Clarke 广义梯度\cite{clarke1976maximum}；分布层面的推广（最优传输/Wasserstein、Schrödinger 桥）超出本章主线，详见规划部扩散章。

\section{本章脉络与展望}
\label{sec:pmp-outlook}

本章把一条三百年的主线严格化：从 1696 年最速降线催生变分法，经 Euler--Lagrange 的必要条件、Legendre--Jacobi--Weierstrass 的二阶充分条件、Hamilton--Noether 的对称与守恒，到 1956 年 Pontryagin 团队用针状变分 + 凸锥分离突破控制约束。下图给出全景。

\begin{figure}[ht]
\centering
\begin{tikzpicture}[>=Stealth, node distance=6mm and 12mm,
  blk/.style={draw, rounded corners, align=center, font=\small, inner sep=3pt, fill=main!6, draw=main!60!black},
  grp/.style={draw, rounded corners, align=center, font=\small, inner sep=3pt, fill=second!8, draw=second!60!black},
  ln/.style={->, thick, gray!55!black}]
  \node[blk] (var) {变分法\\$\delta J=0$};
  \node[blk, right=of var] (el) {Euler--Lagrange\\（无约束）};
  \node[grp, right=of el] (pmp) {PMP\\（有约束·必要）};
  \node[blk, above=8mm of pmp] (ham) {Hamilton 原理\\辛 / Noether};
  \node[grp, below right=8mm and 6mm of pmp] (hjb) {HJB\\（充分·场）};
  \node[grp, below left=8mm and 6mm of pmp] (lqr) {LQR / Riccati\\（解析特例）};
  \node[blk, right=of pmp, xshift=4mm] (ddp) {离散 PMP\\= DDP / iLQR};
  \draw[ln] (var) -- (el);
  \draw[ln] (el) -- (pmp) node[midway, above, font=\scriptsize]{$\max_{\mathbf u}H$};
  \draw[ln] (ham) -- (pmp);
  \draw[ln] (pmp) -- (hjb) node[midway, right, font=\scriptsize]{$\boldsymbol\lambda=\nabla V$};
  \draw[ln] (pmp) -- (lqr);
  \draw[ln] (pmp) -- (ddp);
\end{tikzpicture}
\caption{从变分法到现代最优控制的全景：EL 处理无约束，PMP 以 $\max_{\mathbf u}H$ 突破控制约束，向下归约出 LQR/HJB、向右离散为 DDP/iLQR。}
\label{fig:pmp-panorama}
\end{figure}

这条线索是后续整部最优控制子部的总入口：动态规划与 Bellman/HJB/黏性解的主体在 \cref{ch:dp-hjb}（本章 \cref{eq:pmp-lambda-V} 是其入口）；LQR/LQG/Riccati 的深化在 \cref{ch:lqr-lqg-riccati}（本章给「PMP $\Rightarrow$ Riccati」一步）；Lyapunov 稳定性与共轭点的稳定性互引见 \cref{ch:lyapunov-stability}；鲁棒/$H_\infty$（状态空间 Riccati/博弈视角）见 \cref{ch:robust-hinf}；CLF/CBF 安全与状态约束的相邻见 \cref{ch:clf-cbf}；MPC 深化（间接 vs 直接、PMP $\leftrightarrow$ KKT）见 \cref{ch:mpc-advanced}；DDP/iLQR（离散 PMP $\leftrightarrow$ backward pass）见 \cref{ch:ddp-ilqr}；HJ 可达性见 \cref{ch:hj-reachability}。向动力学部，本章的变分来源 \cref{eq:pmp-robot-dyn}、辛/Noether 接 \cref{ch:spatial-dynamics,ch:geometric-mechanics}；混合 PMP 接 \cref{ch:contact-mechanics,ch:force-wbc}。

\paragraph{历史时间线（订正版）。}
\begin{center}\small
\begin{tabular}{lll}
\toprule
年代 & 人物 & 里程碑\\
\midrule
1696 & Johann Bernoulli & 最速降线挑战，变分法诞生契机\\
1744 & Euler & 系统化 EL 方程\\
1755 & Lagrange & $\delta$ 算法（纯分析方法）；Euler 主动推迟发表己法以让 Lagrange 优先\\
1786/1837/1879 & Legendre/Jacobi/Weierstrass & 二阶充分条件大厦\\
1834 & Hamilton & 正则方程、辛结构\\
1918 & Noether & 对称性 $\leftrightarrow$ 守恒量\\
1956 & Pontryagin 等 & 极大值原理（动机：火箭终端速度最大化，Mayer 形）；首次宣告 Dokl.\\
1958 & Boltyanskii & 一般非线性情形首个严格证明（针状变分 + Brouwer 不动点）\\
1957/1962 & Bellman / Pontryagin 等 & 动态规划 / 专著英译\\
\bottomrule
\end{tabular}
\end{center}

\begin{note}[史料订正]
（1）Lagrange 1755 信引入 $\delta$ 算法属实；Euler 1756 即采纳并把学科改名「calculus of variations」，出于让 Lagrange 优先\emph{推迟}发表自己的讲义（约十年量级，非「将近 30 年」）。（2）PMP 的原始驱动问题是「最大化火箭/导弹终端速度」（Mayer 形），这也是「极大化」约定的由来；「国土防空拦截」是可信但被加工的叙事，非标准史料措辞。（3）非线性一般情形首个严格证明主流归 Boltyanskii（1958，针状变分 + Brouwer 不动点）；另有文献认为 Rozonoer（1959）给出更规范的证明，故宜表述为「Boltyanskii 给出首个严格证明」而不绝对化。
\end{note}

\section*{本章常见误解汇总}
\begin{center}\small
\begin{tabular}{p{0.42\linewidth}p{0.5\linewidth}}
\toprule
误解 & 正确理解\\
\midrule
EL 方程的解一定是极小 & EL 只是一阶必要条件，解可能是极小/极大/鞍点，须二阶条件（Legendre/Jacobi/Weierstrass）验证\\
最小作用量原理说作用量总取最小 & 准确说法是取\emph{驻值}；越过共轭点后实际轨迹可能是鞍点\\
PMP 是充分条件，满足即最优 & PMP 是必要条件，满足者称\emph{极值轨迹}，未必最优，须 HJB 充分性或二阶条件验证\\
针状变分只是一种数学技巧 & 它绕过经典变分在控制约束边界「无法扰动」的本质困难，单点跳变后仍留在 $U$ 内，使非凸/离散 $U$ 也可处理\\
Bang--Bang 是「粗暴」控制 & 它是 $H$ 对 $\mathbf u$ 线性时的精确最优，非近似\\
所有最短时间问题都是 Bang--Bang & 仅当 $H$ 对 $\mathbf u$ 严格线性；$\mathbf u$ 非线性进入 $\mathbf f$ 时可为连续控制\\
共态 $\boldsymbol\lambda$ 只是数学工具 & $\boldsymbol\lambda=\nabla_{\mathbf x}V$（极小化约定），是值函数对状态的梯度（边际代价/影子价格）\\
两约定混用不影响结果 & 极大化与极小化共态差一个负号；接口处必须换算 $\boldsymbol\lambda\to-\boldsymbol\lambda$，否则反馈律差正负号\\
\bottomrule
\end{tabular}
\end{center}

\section*{本章小结}
\begin{center}\small
\begin{longtable}{p{0.26\linewidth}p{0.46\linewidth}p{0.2\linewidth}}
\toprule
知识点 & 核心内容 & 关联\\
\midrule
第一变分 & $\delta J[y;h]=\frac{\mathrm d}{\mathrm d\varepsilon}|_0 J[y+\varepsilon h]$；弱/强变分 & \cref{eq:pmp-gateaux}\\
Euler--Lagrange & $L_y-\frac{\mathrm d}{\mathrm dx}L_{y'}=0$；变分基本引理 & \cref{eq:pmp-el}\\
Beltrami / Noether & $L-y'L_{y'}=C$；连续对称 $\leftrightarrow$ 守恒量 & \cref{thm:pmp-beltrami,eq:pmp-noether}\\
约束变分 & 等周(常数 $\lambda$)/holonomic(函数 $\lambda(x)$)/横截条件 & \cref{eq:pmp-holonomic}\\
二阶条件 & Legendre $L_{y'y'}\ge0$、Jacobi 无共轭点、Weierstrass $\mathcal E\ge0$ & \cref{eq:pmp-jacobi,eq:pmp-weierstrass}\\
PMP 五件套 & 状态/共态/极大化/横截/$H$ 守恒；针状变分证明 & \cref{thm:pmp-main}\\
PMP $\Rightarrow$ EL/HJB/Riccati & $\boldsymbol\lambda$=广义动量；$\boldsymbol\lambda=\nabla V$；Riccati & \cref{eq:pmp-lambda-V,eq:pmp-riccati}\\
极大化结构 & Bang--Bang(线性)/内点(凹)/奇异弧(退化)/Bang-off-Bang & \cref{eq:pmp-switching,eq:pmp-kelley}\\
经典实例 & 双积分器/LQR/Dubins 六型/Goddard 三段式/四旋翼 & \cref{thm:pmp-dubins}\\
现代扩展 & 状态约束(共态跳跃)/离散 PMP/混合系统/$\SE(3)$/Mangasarian & \cref{eq:pmp-discrete,eq:pmp-state-constraint}\\
\bottomrule
\end{longtable}
\end{center}

\noindent\textbf{符号速查.}\;$J$（代价泛函）、$L$（Lagrangian）、$H$（控制 Hamiltonian，极大化约定 $\boldsymbol\lambda^\top\mathbf f-\lambda_0 L$）、$\boldsymbol\lambda$（共态，$\boldsymbol\lambda_{\max}=-\nabla_{\mathbf x}V$）、$\lambda_0$（异常乘子）、$V$（值函数）、$\boldsymbol\sigma$（切换函数）、$\boldsymbol\Phi$（状态转移矩阵）、$\boldsymbol\Omega$（辛矩阵）、$\xi\approx0.4462$（Fuller 切换曲线系数）、$1/\alpha\approx0.242$（Fuller 时间公比）。

\section*{延伸阅读}
\begin{itemize}
  \item \textbf{单本贯通 CoV$\to$PMP$\to$HJB}：Liberzon, \emph{Calculus of Variations and Optimal Control Theory}, Princeton 2012\cite{liberzon2012calculus}——本章主参考，最清晰的现代入门。
  \item \textbf{变分法主体}：Gelfand \& Fomin, \emph{Calculus of Variations}, Dover 1963\cite{gelfand1963calculus}（EL、二阶条件）；Lanczos, \emph{The Variational Principles of Mechanics}\cite{lanczos1970variational}（力学变分的物理直觉）。
  \item \textbf{PMP 原著与史}：Pontryagin 等, \emph{The Mathematical Theory of Optimal Processes}, 1962\cite{pontryagin1962mathematical}；Sussmann \& Willems, ``300 Years of Optimal Control'', \emph{IEEE CSM} 1997\cite{sussmann1997threehundred}（历史综述）。
  \item \textbf{工程视角与例题}：Kirk, \emph{Optimal Control Theory}\cite{kirk2004optimal}；Bryson \& Ho, \emph{Applied Optimal Control}, 1975\cite{bryson1975applied}（航天导向）。
  \item \textbf{专题}：Hartl--Sethi--Vickson 状态约束综述\cite{hartl1995survey}；Anderson \& Moore LQR\cite{anderson2007optimal}；Athans \& Falb Bang--Bang 综合\cite{athans1966optimal}；离散变分力学 Marsden \& West\cite{marsden2001discrete}。
\end{itemize}

本章为整部最优控制理论奠定了一阶必要条件的统一语法。下一章（\cref{ch:dp-hjb}）将从「单条轨迹」转向「全状态空间的值函数」，给出充分条件与最优反馈律——两条路殊途同归于 $\boldsymbol\lambda(t)=\nabla_{\mathbf x}V(\mathbf x^\ast(t),t)$。理解这层对偶，是统一后续所有算法（iLQR、MPC、强化学习）的理论前提。
